Text-to-image models can synthesize detailed scenes from natural language
\citep{ldm,imagen}, with further gains in semantic control and visual fidelity
\citep{dalle2,sdxl}. Safety mechanisms suppress unwanted content during
sampling \citep{sld}, remove concepts from model weights \citep{esd,forgetmenot,uce},
and classify image-policy violations \citep{llavaguard,shieldgemma2}.
An attack that bypasses these mechanisms may still change the requested
content. A nude sculpture, a partly exposed illustration, and an explicit
figure in a specified pose can satisfy the same broad unsafe-content criterion
while demonstrating different capabilities. We study whether a generated
figure jointly retains its intended material, anatomical detail, and pose.

\textbf{Existing methods leave a gap on GPT-Image-2.} We could not locate
public outputs from the compared methods that allow these three requirements
to be verified together on GPT-Image-2 \citep{past2harm,rise}.
Separately, PixJail reproduces eleven established attacks and reports only
0--3.9\% ASR on GPT-Image-2 \citep{pixjail}. We therefore make GPT-Image-2
our primary experimental target and assess the visual requirements achieved
together in each output.

Prior attacks explore token search \citep{sneakyprompt,ringabell}, multimodal
perturbations \citep{mmadiffusion}, and semantic reformulation
\citep{daca,mja}. Low-Effort studies artistic reframing and material
substitution \citep{aestheticcamouflage}; MODX searches artistic modifiers
\citep{modx}. Such changes can also alter the depicted body. Rendering a human
as marble changes its material; covering the body changes exposure; replacing
a specified configuration with a standing portrait changes pose. Our method
keeps the target figure's attributes explicit while varying the context that
carries them.

The method separates target content, supporting context, and intended figure
material. Contextual embedding specifies the relation between the target and
the task represented by its surroundings. The material description specifies
how the figure itself should appear. A photographic outer scene can therefore
contain a drawn human, as in a photographed manga page. Our outputs jointly
attain modern drawn material, maximum exposure under our criteria, and
sexualized posture across different requested configurations.

Material and context ablations examine which changes preserve the target.
Concept-removal evaluation likewise considers whether a system retains
requested content alongside its safety objective \citep{sixcd}. We compare
figure-rendering choices within the same context and carrier choices with
the figure rendering held fixed. Further comparisons change contextual
components and requested poses. These experiments distinguish image return
from preservation of the visual properties that motivated the request.

Evaluating this capability requires more detail than aggregate ASR.
Measurement studies establish the importance of aligned success definitions
and attack configurations \citep{comparisonrequires,asrisnotanumber,singleconfig}.
Safety benchmarks organize risks into categories \citep{t2isafety,t2irisky},
while fine-grained generation benchmarks assess attributes and relations
\citep{tifa,geneval,compbench}. We independently assess the figure's material,
anatomical exposure, and pose, then test their joint agreement with the target.
The resulting profile explains which requirements each output fulfils.
Reporting joint attainment alongside its components connects attack success
to the specific visual capability being evaluated.

\paragraph{Contributions.}
\textbf{(i) A contextual embedding attack.} A prompt-only black-box method
that relates target content to a supporting context while preserving explicit
requirements for the figure's material, exposure, and pose.
\textbf{(ii) Material and context ablations.} Comparisons of figure rendering,
external carriers, and contextual components, supported by outputs that
jointly attain modern drawn material, maximum exposure, and sexualized posture.
\textbf{(iii) A decoupled evaluation framework.} Independent assessment of
material, exposure, and pose, combined with target matching to identify the
specific capability behind an attack's success rate.

